\section{Etapa 9 --- Art\'iculo URUCON y caracterizaci\'on medida de la cadena}
\label{sec:et9}
\periodo{15 --- 23 de julio de 2026 \quad$\cdot$\quad 5 d\'ias de trabajo}

Dos l\'ineas: escribir y someter a revisi\'on el art\'iculo sobre la
auto-calibraci\'on, y ---en paralelo--- \textbf{medir de verdad} la cadena
anal\'ogica con generador y osciloscopio, en vez de confiar en el modelo.

\begin{figure}[ht]
\centering
\includegraphics[width=\textwidth]{identificacion_compensador_medida.png}
\caption{Identificaci\'on experimental hist\'orica del compensador. Los puntos medidos
documentan la respuesta observada; la curva de modelo superpuesta es provisional porque
el primer modelo de operacional confundi\'o GBW con polo dominante.}
\end{figure}

\begin{figure}[ht]
\centering
\includegraphics[width=\textwidth]{geofono_vs_cadena_medida.png}
\caption{Comparaci\'on hist\'orica entre respuesta del ge\'ofono y cadena medida. Se
conserva como diagn\'ostico, no como validaci\'on vigente de la banda compensada hasta
repetir la s\'intesis y la medici\'on.}
\end{figure}

\subsection{El art\'iculo: \emph{Automatic DC-Offset Calibration for a Low-Cost
Geophone Front End}}

\begin{center}\small
\begin{tabularx}{\textwidth}{@{}l X@{}}
\toprule
15 jul & Reescritura de las Secciones III y IV con datos experimentales reales (dos placas, barrido de ganancia completo). Se agrega el \emph{flag} de firmware \texttt{CAL\_PI\_FORCE\_MIN\_DEADBAND} para el experimento decisivo sobre el estancamiento del PGA \\
\addlinespace[2pt]
15 jul & \textbf{Revisi\'on editorial simulada}: se somete el propio art\'iculo a un panel de cinco revisores ficticios antes de mandarlo \\
\addlinespace[2pt]
16 jul & Respuesta a esa revisi\'on: siete referencias nuevas, tabla de resultados ampliada a las cuatro etapas (no solo LP), ecuaciones de banda muerta con justificaci\'on expl\'icita, y \textbf{una correcci\'on de honestidad} sobre qu\'e es realmente la condici\'on ``Default'' \\
\addlinespace[2pt]
23 jul & Reuni\'on con el tutor; compresi\'on de 6 a 5 p\'aginas; figura GEO$\to$ADC lado a lado nueva. Panel de revisi\'on simulado da \emph{Minor Revision}, $\approx$65/100 \\
\bottomrule
\end{tabularx}
\end{center}

\seleccionado{Vale registrar el m\'etodo: \textbf{someter el trabajo a una
revisi\'on adversarial antes de exponerlo}. Lo que produjo no fue cosm\'etico
---la ampliaci\'on de la tabla a las cuatro etapas y la correcci\'on sobre la
condici\'on ``Default'' son cambios de fondo que salieron de esa cr\'itica.}

Del mismo per\'iodo es el cierre de una deuda de nomenclatura: el componente
\texttt{GEO\_ADDER} se renombr\'o a \texttt{GEO\_SUM} en todo el firmware
(esquem\'atico, calibraci\'on, tablas por etapa), con recalibraci\'on de la
ganancia del PI ($\times2{,}63$) y banda muerta configurable por etapa.

\subsection{El protocolo de barridos: c\'omo se midi\'o la cadena}

Este es el documento t\'ecnico m\'as riguroso que produjo el proyecto
(\texttt{PROTOCOLO\_BARRIDOS\_CADENA\_ANALOGICA.md}). Registra c\'omo medir la
respuesta en frecuencia de las cuatro etapas con un generador Agilent 33220A y
un osciloscopio Tektronix TDS2004B:

\begin{itemize}[itemsep=0.2em]
\item CH1 $=$ salida PGA, CH2 $=$ BP, CH3 $=$ sumador (compensador), CH4 $=$ LP;
\item \texttt{SYNC} del generador $\to$ \texttt{EXT TRIG} del osciloscopio, de
modo que el flanco de \emph{sync} marque el comienzo del barrido;
\item \textbf{sondas en $10\times$ obligatoriamente}: en $1\times$ la
capacitancia de la sonda es comparable con el $C_f = \SI{177}{\pico\farad}$ del
nodo BP y lo carga de forma apreciable ---o sea, \emph{el instrumento
modificar\'ia el circuito que quiere medir};
\item nueve bandas de barrido logar\'itmico desde
\SIrange{100}{1}{} m\SI{}{\hertz}--\SI{}{\hertz} hasta
\SIrange{20}{200}{\kilo\hertz}, cada una con su base de tiempo;
\item tres repeticiones por banda cuando sea posible;
\item \textbf{control previo de sondas}: conectar los cuatro canales al mismo
seno y verificar que las razones den $1\angle 0^{\circ}$, para detectar una
sonda descompensada \emph{antes} de atribuirle el error al circuito.
\end{itemize}

\dato{\textbf{Precauci\'on documentada sobre la etapa SUM.} En la banda donde
$BP/PGA \approx -1$, las dos corrientes de entrada del sumador se cancelan en
torno al \SI{95}{\percent}. Eso hace que el cociente medido sea
extremadamente sensible al ruido y a los valores reales de los resistores. La
recomendaci\'on registrada es \emph{abrir temporalmente una de las dos ramas}
para identificar la etapa sin la cancelaci\'on, y usar la medici\'on de cuatro
canales solo como validaci\'on en la configuraci\'on final.}

\subsection{Tratamiento de datos: lo que se hizo bien}

El script de an\'alisis (\texttt{analizar\_barridos\_cadena\_analogica.m},
\texttt{AnalisisCircuito/}) incorpora varias decisiones que conviene no perder:

\begin{itemize}[itemsep=0.2em]
\item \textbf{Modelo del operacional, no ideal}: se registraron
$A_0=\SI{90}{\decibel}$, GBW $=\SI{8}{\mega\hertz}$,
$R_{in}=\SI{35}{\mega\ohm}$,
$R_{out}=\SI{20}{\ohm}$, $C_{in}=\SI{18}{\pico\farad}$. El \emph{slew rate}
(\SI{4.3}{\volt\per\micro\second}) se trata como l\'imite \emph{no lineal} de
amplitud y no como un polo, que es lo correcto. Sin embargo, otros scripts
introdujeron el GBW de \SI{8}{\mega\hertz} como si fuera el polo dominante; esa
parte del resultado queda invalidada. La relaci\'on $f_p=f_t/A_0$ ya se corrigi\'o
en el c\'odigo, pero falta repetir la s\'intesis y el contraste.
\item \textbf{Cancelaci\'on de red por m\'inimos cuadrados}, la misma t\'ecnica
que en el ESP: se estima la fundamental de \SI{50}{\hertz} y se restan
\emph{solo} los arm\'onicos dentro del pasabanda. Se excluye del ajuste un
entorno de tres ciclos alrededor de cada cruce del barrido \textbf{para no
borrar el est\'imulo}, y se omite la regresi\'on si el registro no contiene al
menos dos ciclos de red.
\item \textbf{Recuperaci\'on parcial de capturas saturadas}: una saturaci\'on ya
no descarta la captura entera. Se excluye un ciclo alrededor de cada recorte y
los tramos lineales continuos restantes entran como submuestras
independientes. En la tanda calibrada esto recuper\'o 22 combinaciones
captura-etapa de 9 capturas saturadas.
\item \textbf{Deduplicaci\'on por SHA-256}: se detect\'o que \texttt{XD\_ALL0006}
era copia exacta de \texttt{XD\_ALL0005} en los cuatro canales y se elimin\'o.
\item \textbf{Orden ampliado con criterio}: dos polos extra por cada
operacional en el camino, aumentando tambi\'en el grado del numerador para
conservar el grado relativo (BP 4/3, compensador 5/4, LP 4/2, PGA$\to$ADC
11/8).
\end{itemize}

\subsection{Un conjunto de datos completo, descartado}

\descartado{\textbf{\texttt{Osciloscopio\_descartado\_2026-07-20\_BP\_43k}} ---
una tanda completa de barridos, con sus 156 archivos y m\'as de 20 gr\'aficos ya
producidos, qued\'o archivada como descartada. Motivo: la rama BP ten\'ia
instalado \SI{8200}{\ohm} donde el dise\~no requiere \SI{7501.9}{\ohm}, de modo
que la etapa no estaba realizando la compensaci\'on buscada. Se conserva en el
repositorio en vez de borrarse, con el motivo en el nombre de la carpeta.}

\seleccionado{Conservar el conjunto descartado \emph{con el motivo en el
nombre} es la decisi\'on correcta: permite volver a \'el si aparece una duda, y
evita que alguien lo reprocese sin saber por qu\'e no vale.}

\subsection{La tanda calibrada del 21 de julio}

Se rehicieron los barridos con la rama BP corregida a
$\SI{6.8}{\kilo\ohm} + $ potenci\'ometro. Resultados de identificaci\'on:

\begin{center}\small
\begin{tabularx}{\textwidth}{@{}l r r r X@{}}
\toprule
\textbf{Etapa} & \textbf{Pts.} & \textbf{Coh.\ med.} & \textbf{RMS ajuste} & \textbf{Estado} \\
\midrule
BP (CH2/CH1) & 85 & 0{,}968 & \SI{6.49}{\decibel} & Parcial: polo inferior fuera de banda, alta frecuencia no ideal \\
COMP (CH3/CH1) & 71 & 0{,}964 & \SI{0.38}{\decibel} & Identificaci\'on buena; el pot no es estimable de forma confiable \\
LP (CH4/CH3) & 33 & 0{,}945 & \SI{0.60}{\decibel} & Confiable en la banda \'util \\
PGA$\to$ADC (CH4/CH1) & 61 & \textbf{0{,}995} & \SI{0.20}{\decibel} & Confiable; din\'amica subs\'onica no observable \\
\bottomrule
\end{tabularx}
\end{center}

\dato{Se extendi\'o la medici\'on hasta \SI{10}{\milli\hertz} con la captura
\texttt{XD\_ALL0017} (\SIrange{0.01}{0.05}{\hertz}), libre de saturaci\'on, que
aporta 14 puntos por transferencia con coherencia mediana entre
0{,}988 y 0{,}9998.}

\subsubsection{El hallazgo principal: la muesca no llega a profundidad}

\dato{En $f_0 = \SI{10.2046}{\hertz}$, el modelo completo predice
$CH3/CH1 = \SI{-57.58}{\decibel}$. La medici\'on da $\approx\SI{-31.26}{\decibel}$:
la muesca \textbf{es reproducible} ---dos capturas independientes difieren solo
\SI{0.42}{\decibel}, y dos barridos nuevos concuerdan--- pero queda
\textbf{\SI{26.33}{\decibel} menos profunda} que el objetivo.}

La causa est\'a identificada y es geom\'etrica, no de firmware:

\begin{center}\small
\begin{tabularx}{\textwidth}{@{}l r X@{}}
\toprule
Cursor para cancelaci\'on exacta & \SI{632.556}{\ohm} & \multirow{2}{*}{} \\
Cursor objetivo para $\zeta_0 = 0{,}25$ & \SI{634.539}{\ohm} & \\
\textbf{Separaci\'on entre ambos ajustes} & \textbf{\SI{1.98}{\ohm}} & sobre un potenci\'ometro de recorrido \SI{2}{\kilo\ohm} \\
Posici\'on efectiva estimada por profundidad & $\approx\SI{674}{\ohm}$ & unos \SI{39}{\ohm} por encima del objetivo \\
\bottomrule
\end{tabularx}
\end{center}

\seleccionado{\textbf{El ajuste requerido es de \SI{1.98}{\ohm} sobre un
recorrido de \SI{2000}{\ohm}: una parte en mil.} Esa es la raz\'on real por la
que la compensaci\'on no alcanza su profundidad de dise\~no, y es una conclusi\'on
de dise\~no, no un problema de calibraci\'on: \emph{la sensibilidad del ajuste
est\'a mal escalada}. La conclusi\'on pr\'actica que qued\'o escrita es no seguir
con m\'as barridos anchos, sino recalibrar alrededor de \SI{10.2046}{\hertz} con
mayor sensibilidad vertical en CH3 y repetir un barrido estrecho.}

\dato{\textbf{Lo que s\'i funciona:} el LP se comporta como se espera
(\SI{13.19}{}/\SI{13.20}{\decibel} medidos contra \SI{13.98}{\decibel} del
modelo; \SI{0.435}{\decibel} RMS en toda la banda de ajuste
\SIrange{10.6}{984}{\hertz}), y la cadena PGA$\to$ADC se identifica con
\SI{0.20}{\decibel} y \SI{1.16}{\degree} RMS. \textbf{El circuito es
repetible}; lo que no est\'a logrado es la calibraci\'on fina de la cancelaci\'on.}

\subsection{Lectura resuelta de los tres valores de \texorpdfstring{$\zeta$}{zeta}}

Los archivos de resultados distinguen tres cantidades que se hab\'ian venido
mezclando:
\begin{itemize}[itemsep=0.15em,topsep=0.2em]
\item \texttt{zeta0\_objetivo} $= 0{,}25$, descrito como \emph{``numerador de
Hcompensate''}. Es el amortiguamiento del ge\'ofono, que el numerador del
compensador debe reproducir para cancelar sus polos.
\item \texttt{zeta1\_denominador\_BP}: $980{,}03$ en la tanda descartada del
20/07 y $937{,}40$ en la calibrada del 21/07. Es el amortiguamiento
\emph{nuevo}, el que efectivamente extiende la banda.
\item \texttt{zeta\_realizada\_componentes}: \textbf{$83{,}661$} en la tanda
descartada (con la resistencia de \SI{8200}{\ohm} equivocada) y
\textbf{$0{,}250000000000114$} en la calibrada.
\end{itemize}
Por tanto, \textbf{$83{,}661$ no es un objetivo de dise\~no}: es el valor que
el numerador realizaba por error con la resistencia equivocada. La tanda
calibrada lo lleva a $0{,}25$, mientras que $937{,}40$ permanece como
amortiguamiento del denominador BP. Esta lectura queda adoptada como
autoritativa y corrige la interpretaci\'on anterior del borrador de entrega.

\subsection{Balance de la etapa}

\begin{center}\small
\begin{tabularx}{\textwidth}{@{}l X@{}}
\toprule
\textbf{Se qued\'o} & Protocolo de barridos reproducible; identificaci\'on por \texttt{tfest} con orden justificado; modelo de operacional real; cancelaci\'on de red por m\'inimos cuadrados tambi\'en en el an\'alisis; recuperaci\'on parcial de capturas saturadas; deduplicaci\'on por SHA-256; revisi\'on adversarial simulada como paso previo a publicar \\
\addlinespace[3pt]
\textbf{Se descart\'o} & Tanda completa del 20/07 con BP de \SI{43}{\kilo\ohm} y rama de \SI{8200}{\ohm}; nomenclatura \texttt{GEO\_ADDER}; sondas en $1\times$ \\
\addlinespace[3pt]
\textbf{Qued\'o abierto} & La muesca de compensaci\'on est\'a \SI{26.33}{\decibel} por encima de su profundidad de dise\~no; el ajuste de \SI{1.98}{\ohm} sobre \SI{2}{\kilo\ohm} es impracticable como est\'a y sugiere rehacer el divisor con mayor resoluci\'on \\
\bottomrule
\end{tabularx}
\end{center}
